\documentclass[11pt]{article}
\usepackage[margin=1in]{geometry}
\usepackage{amsmath,amsthm,amssymb,mathtools}
\usepackage[T1]{fontenc}
\usepackage[utf8]{inputenc}

\newtheorem{theorem}{Theorem}
\newtheorem{proposition}{Proposition}
\newtheorem{lemma}{Lemma}
\newtheorem{corollary}{Corollary}
\theoremstyle{remark}
\newtheorem{remark}{Remark}

\newcommand{\sig}[2]{\sigma_{#1#2}}
\newcommand{\sigA}{\sig{2}{1}}
\newcommand{\sigB}{\sig{2}{2}}
\newcommand{\sigC}{\sig{2}{3}}
\newcommand{\sigL}{\sig{1}{2}}
\newcommand{\sigR}{\sig{3}{2}}
\newcommand{\RhoR}{\rho_{R1}}
\newcommand{\CR}{\mathcal R_C}

\title{d5\_223: Surfaced-bundle closure of the \(\sigR\) side}
\author{}
\date{2026-03-21}

\begin{document}
\maketitle

\paragraph{Purpose.}
The previous note d5\_221 closed the dangerous centered-left target
\[
\sigC=[4,1,4]/[2,2,1].
\]
The present note closes the remaining centered-right target
\[
\sigR=[4,4,1]/[2,1,2]
\]
on the same surfaced-bundle semantics.  The key point is that, once the
surviving immediate \(R1\)-branch out of
\[
\sigA=[4,1,4]/[1,2,2]
\]
is normalized on the right strip, the dangerous repeated-end option is exactly
\(\sigR\), and on the common left family the actual two-step continuation is
read off directly from the shared left word \([4,1,4]\).

\paragraph{Setup.}
Fix a modulus \(m>2\), color \(0\), and the standard normalization \(w_0=s_0=0\).
Write \(\partial_i\) for the coordinate step adding \(+1\) in direction \(i\)
modulo \(m\).  In the surfaced common module these coordinate steps commute.
For the centered-right strip
\[
\CR=\{\sigL,\sigB,\sigR\}
=\bigl\{[1,4,4]/[2,1,2],\ [4,1,4]/[2,1,2],\ [4,4,1]/[2,1,2]\bigr\},
\]
consider the common left source family
\[
S_L=\{x:\ \mathrm{layer}(x)=1,\ q(x)=m-1,\ w(x)=m-1,\ s(x)\neq 0\}.
\]
This is exactly the left source family produced by
\texttt{\_build\_candidate} at \(w_0=s_0=0\).

\begin{proposition}[left-side first and second steps on the common family]
\label{prop:sigma32-left-two-step-formulas}
For every \(x\in S_L\), the first two left-side steps of the relevant prototypes
are:
\begin{align*}
\sigA:&\quad x\mapsto \partial_4(x)\mapsto \partial_1\partial_4(x),\\
\sigL:&\quad x\mapsto \partial_1(x)\mapsto \partial_4\partial_1(x),\\
\sigB:&\quad x\mapsto \partial_4(x)\mapsto \partial_1\partial_4(x),\\
\sigR:&\quad x\mapsto \partial_4(x)\mapsto \partial_4^2(x).
\end{align*}
In particular, the actual left-side continuation of \(\sigA\) on \(S_L\)
agrees immediately with \(\sigB\), differs immediately from \(\sigL\), and has
second step different from \(\sigR\).
\end{proposition}

\begin{proof}
This is the direct meaning of the left words in the surfaced builder
\texttt{\_build\_candidate}.  The four relevant left words are
\[
\sigA:\ [4,1,4],\qquad
\sigL:\ [1,4,4],\qquad
\sigB:\ [4,1,4],\qquad
\sigR:\ [4,4,1].
\]
So on the common left source family \(S_L\), the first left step is routed by
those first letters, and the second left step is routed by the second letters.
That gives the displayed formulas.  The last sentence is then immediate from
those formulas.
\end{proof}

\begin{theorem}[\(\sigR\)-closure for every \(m>2\)]
\label{thm:sigma32-closed-all-m-gt-2}
For every modulus \(m>2\), the normalized surviving immediate \(R1\)-image out of
\(\sigA\) does not realize the dangerous centered-right target
\[
\sigR=[4,4,1]/[2,1,2].
\]
More precisely, on the common left family \(S_L\), its actual two-step
left continuation agrees with the centered core \(\sigB=[4,1,4]/[2,1,2]\)
and differs from the repeated-end target \(\sigR\).
\end{theorem}

\begin{proof}
By Proposition~\ref{prop:sigma32-left-two-step-formulas}, for every \(x\in S_L\)
one has
\[
\sigA^2(x)=\partial_1\partial_4(x),
\qquad
\sigB^2(x)=\partial_1\partial_4(x),
\qquad
\sigR^2(x)=\partial_4^2(x).
\]
Thus the actual left-side two-step continuation of \(\sigA\) agrees with
\(\sigB\).

To distinguish it from \(\sigR\), note that every \(x\in S_L\) satisfies
\(q(x)=m-1\).  Therefore
\[
\partial_1\partial_4(x): q=0,
\qquad
\partial_4^2(x): q=m-1.
\]
These are distinct for every \(m>2\).  Hence the normalized surviving immediate
\(R1\)-image out of \(\sigA\) cannot realize \(\sigR\).
\end{proof}

\begin{remark}[the centered core and the companion are not threatened here]
The same proposition also shows that on the common left family the actual first
left step of \(\sigA\) differs from the \(\sigL=[1,4,4]/[2,1,2]\) prototype,
so this note is not merely excluding \(\sigR\) by a weak second-step accident.
At second-step level, \(\sigL\) and \(\sigB\) agree because
\(\partial_4\partial_1=\partial_1\partial_4\), but the live raw burden after
\(\sigC\)-closure was specifically the dangerous repeated-end target \(\sigR\),
so that coincidence causes no problem.
\end{remark}

\begin{corollary}[the remaining raw target in \(T1\) is closed]
Assume the background reduction from d5\_221, namely that after
\(\sigC\)-closure the only live raw target inside \(T1\) is the normalized
centered-right target \(\sigR\).  Then that remaining raw target is closed.
Consequently, under the same background assumptions \((T0,T2,T3)\), the theorem
object \(T1\) is now reduced to background cells rather than a live new local
comparison.
\end{corollary}

\begin{proof}
The theorem above closes the only remaining live raw target identified in
that reduction.
\end{proof}

\paragraph{Checked-range certificate.}
The companion computation in d5\_222 reruns the surfaced bundle checker on
\(m=3,5,7,9,11,13,15,17\).  On that checked range the common left family has
source count \(m(m-1)\), symmetry-reduced microcase count \(1\),
relative immediate displacement
\[
(0,0,0,0,1),
\]
relative actual second displacement
\[
(0,1,0,0,1),
\]
and relative \(\sigR\)-second displacement
\[
(0,0,0,0,2).
\]
So the surfaced symbolic proof above is exactly consistent with the exact
checked bundle computation.

\end{document}
